\section{Momentos de inércia}
\label{sec:dinAcoplada_momentos_inercia}

A matriz de inércia generalizada $M(\vec y\,)$ é introduzida a partir da
forma quadrática da energia cinética, conforme \eqref{eq_T_quadratica_Mq1} e \eqref{eq_M_total_formula}, e em linha com \cite{spong2006robot,featherstone2008}:

\begin{equation}
T =
\frac{1}{2}\,\dot{\vec y}^{\,T}\,M(\vec y\,)\,\dot{\vec y}.
\label{eq_T_quadratica_Mq_reprise}
\end{equation}


Os elementos $m_{kj}(\vec y\,)$ de $M(\vec y\,)$ têm interpretação direta:
os termos da diagonal representam as inércias equivalentes associadas a
cada grau de liberdade, enquanto os termos fora da diagonal expressam os
acoplamentos inerciais entre a atitude da espaçonave e as juntas do
manipulador \cite{Geradin2004}.
Adota-se o vetor de coordenadas generalizadas conforme
\eqref{eq_q_generalizado_espaconave_mr}, isto é:

\begin{equation}
\vec y
=
\begin{bmatrix}
\vec\eta_B \\
\vec\xi
\end{bmatrix}.
\end{equation}

Em que $\vec\eta_B$ descreve a atitude da espaçonave e $\vec\xi$ reúne as
coordenadas articulares do \gls{mr}. Com essa escolha, a matriz
$M(\vec y\,)$ assume a forma particionada usual para sistemas de espaçonave
flutuante \cite{featherstone2008,Wilde2018}:

\begin{equation}
M(\vec y\,) =
\begin{bmatrix}
M_{BB}(\vec y\,) & M_{BM}(\vec y\,) \\
M_{MB}(\vec y\,) & M_{MM}(\vec y\,)
\end{bmatrix}.
\label{eq_M_blocos}
\end{equation}

Em que $M_{BB}$ agrupa as inércias equivalentes associadas à atitude da
espaçonave, $M_{MM}$ contém as contribuições do manipulador e os blocos
$M_{BM}$, $M_{MB}$ representam os termos de acoplamento entre a
espaçonave e o \gls{mr}.

Do ponto de vista construtivo, a expressão
\eqref{eq_M_total_formula} mostra que $M(\vec y\,)$ é obtida
somando-se as contribuições espaciais de cada elo e da espaçonave, isto é:

\begin{equation}
M(\vec y\,) =
\sum_{i=1}^{n}
J_{\text{sp},i}^{T}(\vec y\,)\,
M_{i,\text{spt}}\,
J_{\text{sp},i}(\vec y\,)
+
M_{\text{espaç},\text{spt}}.
\label{eq_M_total_formula}
\end{equation}


Em que $J_{\text{sp},i}(\vec y\,)$ é o jacobiano espacial do elo $i$ e
$M_{i,\text{spt}}$ é o operador de inércia espacial associado à
massa e ao matriz de inércia do elo no formalismo 6D
\cite{featherstone2008}.
Essa construção garante a consistência entre a
cinemática diferencial (via jacobianos espaciais) e a estrutura da
matriz de inércia generalizada empregada nas equações de Lagrange em
\eqref{eq_dinamica_geral}.
